Generative modelling is a cornerstone approach of modern deep learning for modelling the distribution of a complex random variable from samples.
It has seen applications in a wide range of settings, including the modelling of images, text, and video.
Orthogonally, the field of geometric deep learning has developed to incorporate prior knowledge about the geometric structure of a problem into the solutions deep learning systems learn.
The intersection of these domains has proved highly applicable to solving scientific modelling problems with deep learning.

In this thesis we study a particular class of generative models, known as diffusion models, or score-based generative models.
Typical diffusion models model densities supported on a Euclidean space.
We generalise this to a series of settings in Riemannian geometry.

Firstly, we extend the continuous-time diffusion modelling framework to model densities supported on Riemannian manifolds.

Secondly, we extend it to model densities supported on Riemannian manifolds with boundaries.
As a means to accomplishing this we also introduce a new discretisation technique for reflected stochastic differential equations.

Finally, we extend it to model densities supported on the spaces of tensor fields on Riemannian manifolds, and on the spaces of paths on Riemannian manifolds.
In tandem with this we introduce new techniques for conditionally sampling from diffusion models.

In total these contributions allow for the application of diffusion models to a wide range of scientific problems, and we demonstrate this applicability in numerous settings.